% !TeX program = lualatex
% Copyright (C) 2026 IEEE UFG Computer Society contributors.
% Distributed under the LaTeX Project Public License 1.3c or later.
% See LICENSE and development/MANIFEST.md.
%
% PLANO DE EXECUÇÃO DE PROJETO
% Apresente à Diretoria antes de iniciar o projeto. Após a decisão, registre
% versão, data e referência da ata. Atualize o histórico ao mudar o escopo.
% Substitua todos os campos entre colchetes e remova quadros desnecessários.
\documentclass[aspectratio=169,11pt]{beamer}
\usetheme[sectionpages=false]{ieeeufgcs}
\usepackage[brazilian]{babel}
\usepackage{tabularx}

% Para citar referências, renomeie references.bib.example para
% references.bib e ative estas linhas e o quadro ao final.
% \usepackage{csquotes}
% \usepackage[backend=biber,style=numeric]{biblatex}
% \addbibresource{references.bib}

\newcommand{\ProjetoNome}{[Nome do projeto]}
\newcommand{\ProjetoCodigo}{[Código]}
\newcommand{\ProjetoLider}{[Nome da liderança discente]}
\newcommand{\PlanoVersao}{0.1}
\newcommand{\PlanoData}{[data]}

\title[Plano de projeto]{Plano de execução: \ProjetoNome}
\subtitle{Proposta para decisão da Diretoria Executiva}
\author{Diretoria de Pesquisa e Desenvolvimento e \ProjetoLider}
\institute{Capítulo Estudantil da IEEE Computer Society na UFG}
\date{Versão \PlanoVersao{} · \PlanoData}

\begin{document}

\begin{frame}[plain,noframenumbering]
  \titlepage
\end{frame}

\begin{frame}{Identificação e decisão solicitada}
  \small
  \begin{tabularx}{\textwidth}{
    >{\raggedright\arraybackslash}p{.28\textwidth}
    >{\raggedright\arraybackslash}X
  }
    \textbf{Projeto} & \ProjetoNome{} (\ProjetoCodigo) \\
    \textbf{Frente} & [Hardware e Estruturas ou Software] \\
    \textbf{Liderança} & \ProjetoLider \\
    \textbf{Equipe} & [Três a seis integrantes; nomes ou referência] \\
    \textbf{Situação} & [Proposta, aprovada ou em revisão] \\
    \textbf{Versão e data} & \PlanoVersao{} · \PlanoData \\
  \end{tabularx}
  \vfill
  \begin{block}{Pedido à Diretoria}
    [Aprovar o escopo, os responsáveis, o cronograma e o teto de recursos
    descritos nesta versão, com as condições indicadas.]
  \end{block}
\end{frame}

\begin{frame}{Necessidade, proposta e sucesso}
  \begin{columns}[T,onlytextwidth]
    \begin{column}{.31\textwidth}
      {\IEEECSHeadingFont\bfseries\color{IEEEBlue}Necessidade\par}
      \medskip
      [Problema observável, público afetado e evidência da necessidade.]
    \end{column}
    \begin{column}{.31\textwidth}
      {\IEEECSHeadingFont\bfseries\color{IEEEBlue}Proposta\par}
      \medskip
      [Solução inicial e resultado útil a entregar.]
    \end{column}
    \begin{column}{.31\textwidth}
      {\IEEECSHeadingFont\bfseries\color{IEEEBlue}Sucesso\par}
      \medskip
      [Condição mensurável para aceitar o resultado no prazo.]
    \end{column}
  \end{columns}
  \vfill
  \textbf{Vínculo com o Capítulo:}
  [Formação, extensão, pesquisa ou produção científica esperada.]
\end{frame}

\begin{frame}{Escopo e condições de execução}
  \begin{columns}[T,onlytextwidth]
    \begin{column}{.47\textwidth}
      {\IEEECSHeadingFont\bfseries\color{IEEEBlue}Incluído\par}
      \begin{itemize}
        \item \emph{[Funcionalidade ou atividade essencial]}
        \item \emph{[Documentação e demonstração previstas]}
      \end{itemize}
      \medskip
      {\IEEECSHeadingFont\bfseries\color{IEEEBlue}Fora do escopo\par}
      \begin{itemize}
        \item \emph{[Entrega adiada ou excluída nesta versão]}
      \end{itemize}
    \end{column}
    \begin{column}{.47\textwidth}
      {\IEEECSHeadingFont\bfseries\color{IEEEBlue}Condições\par}
      \begin{itemize}
        \item \emph{[Recurso, espaço ou autorização necessária]}
        \item \emph{[Licenças, dados e segurança a observar]}
        \item \emph{[Limite de tempo da equipe e do ciclo]}
      \end{itemize}
    \end{column}
  \end{columns}
\end{frame}

\begin{frame}{Entregas e critérios de aceitação}
  \small
  \begin{tabularx}{\textwidth}{
    >{\raggedright\arraybackslash}p{.20\textwidth}
    >{\raggedright\arraybackslash}X
    >{\raggedright\arraybackslash}p{.28\textwidth}
  }
    \textbf{Entrega} & \textbf{Critério verificável} &
    \textbf{Evidência} \\
    \hline
    D1: [protótipo] & [Teste ou demonstração que deve passar] &
    [Registro ou link] \\
    D2: [docs] & [Conteúdo mínimo revisado] &
    [Arquivo versionado] \\
    D3: [apresentação] & [Público e resultado esperado] &
    [Evento ou relatório] \\
  \end{tabularx}
  \vfill
  \begin{block}{Condição de conclusão}
    [Quem avalia as entregas, quando e como registra o aceite.]
  \end{block}
\end{frame}

\begin{frame}{Equipe e responsabilidades}
  \small
  \begin{tabularx}{\textwidth}{
    >{\raggedright\arraybackslash}p{.26\textwidth}
    >{\raggedright\arraybackslash}p{.25\textwidth}
    >{\raggedright\arraybackslash}X
  }
    \textbf{Função} & \textbf{Pessoa} & \textbf{Responsabilidade} \\
    \hline
    Liderança discente & \ProjetoLider &
    [Coordenação, apresentações e decisões da equipe] \\
    Equipe técnica & [Integrantes] &
    [Construção, testes e documentação] \\
    Diretoria de P\&D & [Nome] &
    [Plano, acompanhamento e relatório] \\
    Orientação docente & [Nomes] &
    [Supervisão técnica e institucional] \\
  \end{tabularx}
  \vfill
  \textbf{Ritmo de trabalho:} [Reuniões da equipe e apresentação mensal
  à Diretoria.]
\end{frame}

\begin{frame}{Cronograma e marcos}
  \small
  \begin{tabularx}{\textwidth}{
    >{\raggedright\arraybackslash}p{.16\textwidth}
    >{\raggedright\arraybackslash}p{.22\textwidth}
    >{\raggedright\arraybackslash}X
    >{\raggedright\arraybackslash}p{.23\textwidth}
  }
    \textbf{Etapa} & \textbf{Período} & \textbf{Marco} &
    \textbf{Evidência} \\
    \hline
    Preparação & [Datas] & [Equipe e recursos confirmados] &
    [Aprovação e registro] \\
    Primeira versão & [Datas] & [Protótipo mínimo testado] &
    [Teste ou demonstração] \\
    Evolução & [Datas] & [Melhorias priorizadas] &
    [Revisão de resultados] \\
    Encerramento & [Datas] & [Entrega e aceite] &
    [Relatório e repositório] \\
  \end{tabularx}
  \vfill
  \textbf{Dependências do calendário:}
  [Provas, compras, laboratório ou autorizações.]
\end{frame}

\begin{frame}{Recursos e orçamento}
  \small
  \begin{tabularx}{\textwidth}{
    >{\raggedright\arraybackslash}X
    >{\raggedright\arraybackslash}p{.16\textwidth}
    >{\raggedright\arraybackslash}p{.25\textwidth}
    >{\raggedright\arraybackslash}p{.18\textwidth}
  }
    \textbf{Recurso} & \textbf{Qtd.} & \textbf{Fonte/situação} &
    \textbf{Estimativa} \\
    \hline
    \emph{[Material ou serviço]} & [Qtd.] &
    [Confirmado ou pendente] &
    [R\$ valor] \\
    \emph{[Equipamento existente]} & [Qtd.] &
    [Disponível ou a confirmar] &
    [R\$ valor] \\
    \emph{[Espaço ou plataforma]} & [Prazo] &
    [Autorização necessária] &
    [R\$ valor] \\
  \end{tabularx}
  \vfill
  \textbf{Total estimado:} [R\$ valor]. \quad
  \textbf{Fonte confirmada:} [Origem e valor].

  \medskip
  Valores não confirmados exigem nova avaliação antes de gerar despesa.
\end{frame}

\begin{frame}{Riscos e dependências}
  \small
  \begin{tabularx}{\textwidth}{
    >{\raggedright\arraybackslash}p{.27\textwidth}
    >{\raggedright\arraybackslash}p{.21\textwidth}
    >{\raggedright\arraybackslash}X
  }
    \textbf{Risco ou dependência} & \textbf{Sinal de alerta} &
    \textbf{Resposta e responsável} \\
    \hline
    \emph{[Atraso de componentes]} & [Data sem entrega] &
    [Alternativa e pessoa responsável] \\
    \emph{[Sobrecarga da equipe]} & [Marco não cumprido] &
    [Reduzir escopo ou pausar] \\
    \emph{[Risco técnico ou de dados]} & [Teste ou revisão] &
    [Controle, autorização e supervisão] \\
  \end{tabularx}
  \vfill
  \textbf{Condição para iniciar:}
  [Liste dependências que precisam estar resolvidas.]
\end{frame}

\begin{frame}{Acompanhamento e controle de mudanças}
  \begin{columns}[T,onlytextwidth]
    \begin{column}{.47\textwidth}
      {\IEEECSHeadingFont\bfseries\color{IEEEBlue}Acompanhamento\par}
      \begin{itemize}
        \item \emph{[Indicador e fonte dos resultados]}
        \item \emph{[Local dos arquivos e das evidências]}
        \item \emph{[Revisão mensal com a Diretoria]}
      \end{itemize}
    \end{column}
    \begin{column}{.47\textwidth}
      {\IEEECSHeadingFont\bfseries\color{IEEEBlue}Alterações\par}
      \begin{itemize}
        \item \emph{[Quem propõe e justifica a mudança]}
        \item \emph{[Impacto em prazo, custo, risco e entrega]}
        \item \emph{[Decisão, ata e nova versão do plano]}
      \end{itemize}
    \end{column}
  \end{columns}
\end{frame}

\begin{frame}{Histórico de versões e mudanças}
  \small
  \begin{tabularx}{\textwidth}{
    >{\raggedright\arraybackslash}p{.12\textwidth}
    >{\raggedright\arraybackslash}p{.16\textwidth}
    >{\raggedright\arraybackslash}X
    >{\raggedright\arraybackslash}p{.25\textwidth}
  }
    \textbf{Versão} & \textbf{Data} & \textbf{Mudança e impacto} &
    \textbf{Decisão e registro} \\
    \hline
    \PlanoVersao & \PlanoData & Criação da proposta. &
    [Decisão pendente] \\
    \emph{[vX.Y]} & [Data] &
    [Alteração de escopo; efeito em prazo, custo e risco] &
    [Decisão e ata] \\
  \end{tabularx}
  \vfill
  \begin{block}{Ao revisar o plano}
    [Atualize os quadros afetados, registre a justificativa e apresente a
    nova versão à Diretoria para decisão.]
  \end{block}
\end{frame}

\begin{frame}{Aprovação e próximos passos}
  \small
  \begin{tabularx}{\textwidth}{
    >{\raggedright\arraybackslash}p{.27\textwidth}
    >{\raggedright\arraybackslash}X
  }
    \textbf{Versão avaliada} & \PlanoVersao{} · \PlanoData \\
    \textbf{Decisão} & [Aprovado, aprovado com condições ou devolvido] \\
    \textbf{Responsáveis} & [Diretoria e orientação consultada] \\
    \textbf{Registro} & [Ata nº, data e local do arquivo] \\
    \textbf{Condições} & [Pendências antes do início, se houver] \\
  \end{tabularx}
  \vfill
  \begin{block}{Após a aprovação}
    [Confirmar recursos e equipe; publicar a versão aprovada; iniciar o
    primeiro marco e agendar a apresentação mensal.]
  \end{block}
\end{frame}

% Ative este quadro depois de configurar o arquivo references.bib.
% \begin{frame}[allowframebreaks]{Referências}
%   \printbibliography
% \end{frame}

\end{document}
